\subsection{Immutable Set Variants}

Python also provides an immutable set-like object: \texttt{frozenset}. A normal \texttt{set} is mutable. Elements can be added and removed after construction. A \texttt{frozenset} has the same basic membership-domain meaning, but its contents cannot be changed after creation.

This matters because hash-based containers require stable hash behavior. A mutable \texttt{set} cannot be used as a dictionary key or as an element inside another set. A \texttt{frozenset}, however, can be hashable if all of its elements are hashable.

\begin{quote}
A \texttt{set} is a mutable uniqueness domain. A \texttt{frozenset} is a fixed uniqueness domain.
\end{quote}

\subsubsection{\texttt{frozenset} as an Immutable Hashable Set-Like Object}

A \texttt{frozenset} can be created by passing an iterable to the \texttt{frozenset()} constructor.

\begin{verbatim}
names = frozenset(["Arthur", "Brian", "Arthur", "Lancelot"])

print(f"names: {names!r}")
print(f"type(names): {type(names)}")
print(f"len(names): {len(names)}")
\end{verbatim}

Possible output:

\begin{verbatim}
names: frozenset({'Arthur', 'Brian', 'Lancelot'})
type(names): <class 'frozenset'>
len(names): 3
\end{verbatim}

The duplicate value \texttt{"Arthur"} appears only once. A \texttt{frozenset} still obeys the uniqueness rule of sets.

The main difference is mutation. A normal \texttt{set} has methods such as \texttt{add()}, \texttt{remove()}, \texttt{discard()}, \texttt{pop()}, and \texttt{clear()}. A \texttt{frozenset} does not have these mutating methods.

\begin{verbatim}
normal_set = {"spam", "eggs"}
fixed_set = frozenset(["spam", "eggs"])

print("normal_set")
print(f"  type: {type(normal_set)}")
print(f"  add in dir: {'add' in dir(normal_set)}")
print(f"  remove in dir: {'remove' in dir(normal_set)}")
print(f"  union in dir: {'union' in dir(normal_set)}")
print()

print("fixed_set")
print(f"  type: {type(fixed_set)}")
print(f"  add in dir: {'add' in dir(fixed_set)}")
print(f"  remove in dir: {'remove' in dir(fixed_set)}")
print(f"  union in dir: {'union' in dir(fixed_set)}")
\end{verbatim}

Possible output:

\begin{verbatim}
normal_set
  type: <class 'set'>
  add in dir: True
  remove in dir: True
  union in dir: True

fixed_set
  type: <class 'frozenset'>
  add in dir: False
  remove in dir: False
  union in dir: True
\end{verbatim}

The \texttt{frozenset} keeps non-mutating set operations such as membership testing, union, intersection, difference, and symmetric difference. It rejects operations that would change its contents in place.

\begin{verbatim}
fixed_set = frozenset(["spam", "eggs", 42])

print(f"fixed_set: {fixed_set!r}")
print(f"{'spam' in fixed_set = }")
print(f"{'larch' in fixed_set = }")
print(f"{42 in fixed_set = }")
\end{verbatim}

Possible output:

\begin{verbatim}
fixed_set: frozenset({'eggs', 42, 'spam'})
'spam' in fixed_set = True
'larch' in fixed_set = False
42 in fixed_set = True
\end{verbatim}

Trying to mutate a \texttt{frozenset} fails because the method does not exist.

\begin{verbatim}
fixed_set = frozenset(["spam", "eggs"])

try:
    fixed_set.add("larch")
except AttributeError as err:
    print(f"mutation failed: {err}")
\end{verbatim}

Possible output:

\begin{verbatim}
mutation failed: 'frozenset' object has no attribute 'add'
\end{verbatim}

This is not a temporary restriction. It is the defining property of the object. Once created, the membership domain is fixed.

A \texttt{frozenset} is also hashable when all its elements are hashable.

\begin{verbatim}
fixed_set = frozenset(["spam", "eggs", 42])

print(f"fixed_set: {fixed_set!r}")
print(f"type(fixed_set): {type(fixed_set)}")
print(f"hash(fixed_set): {hash(fixed_set)}")
\end{verbatim}

Possible output:

\begin{verbatim}
fixed_set: frozenset({'eggs', 42, 'spam'})
type(fixed_set): <class 'frozenset'>
hash(fixed_set): -7694472856633011256
\end{verbatim}

The exact hash value may differ between Python runs. The important point is that the object can produce a hash value because its membership cannot change after construction.

A mutable \texttt{set} cannot do this.

\begin{verbatim}
normal_set = {"spam", "eggs"}

print(f"normal_set: {normal_set!r}")
print(f"type(normal_set): {type(normal_set)}")
print(f"normal_set.__hash__: {normal_set.__hash__}")

try:
    print(hash(normal_set))
except TypeError as err:
    print(f"hash failed: {err}")
\end{verbatim}

Possible output:

\begin{verbatim}
normal_set: {'eggs', 'spam'}
type(normal_set): <class 'set'>
normal_set.__hash__: None
hash failed: unhashable type: 'set'
\end{verbatim}

The normal set is unhashable because its contents can change. If its hash depended on its current elements, adding or removing an element would change the hash-relevant state. That would be unsafe inside another hash table.

A \texttt{frozenset} can still be iterated.

\begin{verbatim}
fixed_set = frozenset(["Homer", "Marge", "Lisa"])

print(f"fixed_set: {fixed_set!r}")

for name in fixed_set:
    print(f"name: {name!r}")
\end{verbatim}

Possible output:

\begin{verbatim}
fixed_set: frozenset({'Lisa', 'Homer', 'Marge'})
name: 'Lisa'
name: 'Homer'
name: 'Marge'
\end{verbatim}

The iteration order must not be treated as a meaningful stored order. A \texttt{frozenset}, like a \texttt{set}, is still unordered.

\subsubsection{Using \texttt{frozenset} as a Dictionary Key or Set Element}

The practical reason \texttt{frozenset} exists is that sometimes the program needs a set-like value to behave as one stable hashable object.

A dictionary key must be hashable. A normal set cannot be used as a key.

\begin{verbatim}
bad_key = {"spam", "eggs"}

try:
    data = {bad_key: "breakfast"}
except TypeError as err:
    print(f"dictionary construction failed: {err}")
\end{verbatim}

Possible output:

\begin{verbatim}
dictionary construction failed: unhashable type: 'set'
\end{verbatim}

A \texttt{frozenset} can be used instead.

\begin{verbatim}
good_key = frozenset(["spam", "eggs"])

data = {
    good_key: "breakfast"
}

print(f"good_key: {good_key!r}")
print(f"type(good_key): {type(good_key)}")
print(f"data: {data!r}")
print(f"data[good_key]: {data[good_key]!r}")
\end{verbatim}

Possible output:

\begin{verbatim}
good_key: frozenset({'eggs', 'spam'})
type(good_key): <class 'frozenset'>
data: {frozenset({'eggs', 'spam'}): 'breakfast'}
data[good_key]: 'breakfast'
\end{verbatim}

This is useful when the key represents an unordered group rather than a sequence. For example, a pair of people can be treated as the same group regardless of order.

\begin{verbatim}
pair1 = frozenset(["Jerry", "George"])
pair2 = frozenset(["George", "Jerry"])

print(f"pair1: {pair1!r}")
print(f"pair2: {pair2!r}")
print(f"pair1 == pair2: {pair1 == pair2}")
print(f"hash(pair1) == hash(pair2): {hash(pair1) == hash(pair2)}")
\end{verbatim}

Possible output:

\begin{verbatim}
pair1: frozenset({'George', 'Jerry'})
pair2: frozenset({'George', 'Jerry'})
pair1 == pair2: True
hash(pair1) == hash(pair2): True
\end{verbatim}

Because the two \texttt{frozenset} objects contain the same members, they compare equal. Since they compare equal, they must also have equal hash values.

This makes them safe dictionary keys.

\begin{verbatim}
dialogue = {}

pair1 = frozenset(["Jerry", "George"])
pair2 = frozenset(["George", "Jerry"])

dialogue[pair1] = "coffee shop argument"
dialogue[pair2] = "updated coffee shop argument"

print(f"dialogue: {dialogue!r}")
print(f"len(dialogue): {len(dialogue)}")
print(f"dialogue[pair1]: {dialogue[pair1]!r}")
print(f"dialogue[pair2]: {dialogue[pair2]!r}")
\end{verbatim}

Possible output:

\begin{verbatim}
dialogue: {frozenset({'George', 'Jerry'}): 'updated coffee shop argument'}
len(dialogue): 1
dialogue[pair1]: 'updated coffee shop argument'
dialogue[pair2]: 'updated coffee shop argument'
\end{verbatim}

The second assignment did not create a second dictionary entry. It used an equal key and replaced the associated value.

A \texttt{frozenset} can also be stored as an element inside a normal set.

\begin{verbatim}
groups = set()

group1 = frozenset(["Arthur", "Brian"])
group2 = frozenset(["Brian", "Arthur"])
group3 = frozenset(["Homer", "Marge"])

groups.add(group1)
groups.add(group2)
groups.add(group3)

print(f"groups: {groups!r}")
print(f"type(groups): {type(groups)}")
print(f"len(groups): {len(groups)}")
\end{verbatim}

Possible output:

\begin{verbatim}
groups: {frozenset({'Arthur', 'Brian'}), frozenset({'Homer', 'Marge'})}
type(groups): <class 'set'>
len(groups): 2
\end{verbatim}

The first two groups collapse into one set element because they are equal \texttt{frozenset} objects.

A normal set cannot be placed inside another set.

\begin{verbatim}
outer = set()
inner = {"Arthur", "Brian"}

try:
    outer.add(inner)
except TypeError as err:
    print(f"nested set failed: {err}")
\end{verbatim}

Possible output:

\begin{verbatim}
nested set failed: unhashable type: 'set'
\end{verbatim}

The fix is to freeze the inner set-like group.

\begin{verbatim}
outer = set()
inner = frozenset(["Arthur", "Brian"])

outer.add(inner)

print(f"outer: {outer!r}")
print(f"type(outer): {type(outer)}")
\end{verbatim}

Possible output:

\begin{verbatim}
outer: {frozenset({'Arthur', 'Brian'})}
type(outer): <class 'set'>
\end{verbatim}

The practical distinction is:

\begin{center}
\begin{tabular}{lll}
\textbf{Object} & \textbf{Mutable?} & \textbf{Hashable?} \\
\hline
\texttt{set} & Yes & No \\
\texttt{frozenset} & No & Yes, if all elements are hashable \\
\end{tabular}
\end{center}

A \texttt{frozenset} should be used when the program needs an unordered unique group that must itself be stored inside a hash-based structure.

In ordinary words:

\begin{itemize}
    \item use \texttt{set} for a membership domain that must change;
    \item use \texttt{frozenset} for a membership domain that must remain fixed;
    \item use \texttt{frozenset} when a set-like group must become a dictionary key or another set's element.
\end{itemize}